%!TEX root = ../main.tex

\chapter*{Zusammenfassung / Abstract}

Die Applikation \ac{ATLAS} des Laser Application Center dient der Dokumentation und Auswertung von Laser-Schweißversuchen und besteht aus einem Angular-Frontend sowie mehreren ASP.NET-Core-Backend-Services. Zu Beginn dieser Arbeit war die bestehende Testinfrastruktur unvollständig. Mehrere zentrale Backend-Services verfügten über kein eigenes Testprojekt, und die Testabdeckung im Frontend lag unterhalb üblicher Zielwerte. Ziel der Arbeit ist die Weiterentwicklung einer robusten und wartbaren Testinfrastruktur für \ac{ATLAS}, durch die Erstellung neuer Tests für zentrale Komponenten sowie durch ein Konzept für deren automatisierte Ausführung in einer eigenen Testumgebung.

Zur Erreichung dieses Ziels wurden zunächst die bestehende Systemarchitektur und der Ist-Zustand der vorhandenen Tests analysiert. Auf dieser Grundlage wurden funktionale und nicht-funktionale Anforderungen an die Testinfrastruktur abgeleitet und priorisiert. Darauf aufbauend wurde ein zweistufiges Konzept entwickelt, das zunächst eine lokal ausführbare Testumgebung und anschließend eine mögliche zentrale Ausführungsumgebung auf Basis einer virtuellen Maschine vorsieht. Im praktischen Teil der Arbeit wurden für priorisierte Backend-Services Unit- und Integrationstests sowie für zentrale Frontend-Bereiche Unit-Tests und ergänzende UI-End-to-End-Tests mit Playwright umgesetzt.

Die Umsetzung führte zu einer messbaren Verbesserung der Testabdeckung in Frontend und Backend. Im Frontend stieg die Unit-Test-Coverage in allen vier betrachteten Kennzahlen um jeweils drei bis vier Prozentpunkte. Im Backend wurde die Testabdeckung der bearbeiteten Services deutlich gesteigert, für zwei Services entstand dabei erstmals ein eigenes Testprojekt. Die zentrale Testumgebung blieb aus zeitlichen Gründen auf Konzeptebene und wurde nicht produktiv umgesetzt. Die Ergebnisse zeigen, dass sich die Testabdeckung gezielt verbessern lässt. Gleichzeitig bleibt weiterer Ausbau erforderlich, um bestehende Lücken systematisch zu schließen und die Testausführung künftig unabhängig von lokalen Entwicklungsumgebungen zu ermöglichen.

\vspace{1em}

\textit{The application \ac{ATLAS} of the Laser Application Center supports the documentation and evaluation of laser welding experiments and consists of an Angular frontend together with several ASP.NET Core backend services. At the start of this thesis, the existing test infrastructure was incomplete. Several central backend services had no dedicated test project, and frontend test coverage fell below common target values. The aim of this thesis is to further develop a robust and maintainable test infrastructure for \ac{ATLAS}, through the creation of new tests for central components and through a concept for their automated execution in a dedicated test environment.}

\textit{To achieve this goal, the existing system architecture and the current state of the available tests were first analysed. Based on this analysis, functional and non-functional requirements for the test infrastructure were derived and prioritised. Building on these requirements, a two-stage concept was developed, comprising a locally executable test environment followed by a possible central execution environment based on a virtual machine. In the practical part of the thesis, unit and integration tests were implemented for prioritised backend services, and unit tests together with additional UI-end-to-end tests using Playwright were implemented for central frontend areas.}

\textit{The implementation resulted in a measurable improvement of test coverage in both the frontend and the backend. In the frontend, unit test coverage increased by three to four percentage points across all four examined metrics. In the backend, test coverage for the addressed services was substantially increased, with a dedicated test project being established for two services for the first time. Due to time constraints, the central test environment remained at the concept stage and was not implemented in production. The results show that test coverage can be improved in a targeted manner. Further work remains necessary to systematically close existing gaps and to enable test execution independent of local development environments in the future.}